\documentclass{amsart}

\input{../preamble_general.tex}

\author{Jalen Chrysos}
\title{Hamiltonian Paths and Cycles}

\begin{document}
	
	\maketitle 
	
	A \textit{Hamiltonian Path} is a path which includes every vertex exactly once. Similarly, a \textit{Hamiltonian Cycle} is a cycle which includes every vertex exactly once. In general, it is not easy to tell whether a graph has a Hamiltonian path/cycle, unlike with Eulerian paths and cycles, where one can determine the existence of a path by checking vertices' degrees.\\
	
	\noindent \textbf{Problem 1}: Let $G_{a,b}$ be the graph consisting of $ab$ vertices arranged in an $a\times b$ grid, with adjacent vertices connected by an edge. 
	\begin{enumerate}[(i)]
		\item For which $a,b$ does $G_{a,b}$ have a Hamiltonian path?
		\item For which $a,b$ does $G_{a,b}$ have a Hamiltonian cycle?\\
	\end{enumerate}
	
%	\noindent \textbf{Problem 2}: Which platonic solid graphs have a Hamiltonian path / cycle?\\
%	
%	$$\textbf{illustration of the platonic solids}$$
	
	\noindent \textbf{Problem 2}: Show that if $G$ has a Hamiltonian path, then any induced subgraph $G^*$ obtained by deleting $k$ vertices of $G$ must have at most $k+1$ connected components.\\
	
	\noindent \textbf{Problem 3}: An (open) \textit{knight's tour} is a sequence of knight moves on a rectangular grid such that the knight visits every square exactly once (the start and end squares count as visited). Show that that a knight's tour is impossible on the following board sizes:
	\begin{enumerate}[(i)]
		\item $3\times 3$
		\item $4\times 4$
		\item $3\times 6$\\
	\end{enumerate}
	
	\noindent \textbf{Problem 4}: Find a \textit{connected} 3-regular graph (i.e. graph whose vertices all have degree 3) without any Hamiltonian path.\\
	
	\noindent \textbf{Problem 5}: Let $G_n$ be the $n$-dimensional cube graph; that is, $G_n$ has $2^n$ vertices corresponding to binary strings $(0,1,1,0,\dots)$ of length $n$, and two strings have an edge between them if and only if they differ in exactly one place (e.g. $(0,1,1)\sim (1,1,1)$ but $(0,1,1)\not\sim (1,0,1)$). Show that $G_n$ has a Hamiltonian cycle for all $n\geq 2$.\\
	
\noindent \textbf{Problem 6} ($\star$): Let $G$ be a graph on $n\geq 3$ vertices.
\begin{enumerate}[(i)]
	\item Show that if $G$ has no Hamiltonian cycle, and $x,y$ are vertices such that $x\not\sim y$ and
	$$
	\deg(x)+\deg(y)\geq n,
	$$
	then the graph $G^+$ obtained by adding an edge $x\sim y$ also has no Hamiltonian cycle\footnote{If one takes a graph $G$ and repeatedly adds edges between any two vertices $x,y$ for which $\deg(x)+\deg(y)\geq n$, the resulting graph is called the \textit{closure} of $G$. This problem shows that $G$ has a Hamiltonian cycle if and only if its closure does.}.
	\item (Ore's Theorem): Show that if $\deg(x)+\deg(y)\geq n$ for all $x\not\sim y$, $G$ has a Hamiltonian cycle.
	\item (P\'osa's Theorem): Show that if $G$ has no Hamiltonian cycle, then for some $1\leq d<n/2$, there are at least $d$ vertices with degree $\leq d$. (\textit{Hint: It suffices to show this for the closure of $G$. Consider the set of non-neighbors of a suitable vertex $x$}).
	\item Suppose that $n=9$ and $G$ has score $(2,3,4,5,\ast,\ast,\ast,\ast,\ast)$, where each $\ast$ is some number $\geq 5$. Show that $G$ has a Hamiltonian cycle.\\
\end{enumerate}
	
	
\end{document}